\documentclass[a4paper,10pt]{article}
\usepackage[english]{babel}
%Includes "References" in the table of contents
\usepackage[nottoc]{tocbibind}
%%
\usepackage{graphicx}
\graphicspath{{./Figures/}}
\usepackage{enumitem}

%\usepackage{subcaption}
%\captionsetup[table]{font=small,size=smaller,textfont=sc}
%\captionsetup[figure]{font=small,size=smaller}

\usepackage{booktabs}
%\usepackage[T1]{fontenc}
%\usepackage{cite}
%\usepackage[cmex10]{amsmath}
%\usepackage{color}
%\usepackage{bm,bbm}
%\usepackage{wasysym}
%\usepackage{texnames}
%\usepackage{url}

\usepackage[boxed]{algorithm2e}   % AAB inserido
%\usepackage{natbib} % AAB Inserted
\usepackage[listings]{tcolorbox}

%\usepackage{siunitx}
\usepackage{natbib}
\usepackage{multirow,bigstrut}
\usepackage{subfig} % AAB inserted	
\usepackage{booktabs} % AAB inserted
%\usepackage[detect-weight=true, binary-units=true]{siunitx} % AAB inserted
\usepackage{siunitx} % AAB inserted
\usepackage{tikz}  % AAB inserido
\usetikzlibrary{shapes,arrows,shadows}  % AAB inserido
\usepackage{bbm}
%%

%DIF PREAMBLE EXTENSION ADDED BY LATEXDIFF
%DIF UNDERLINE PREAMBLE %DIF PREAMBLE
\RequirePackage[normalem]{ulem} %DIF PREAMBLE
\RequirePackage{color}\definecolor{RED}{rgb}{1,0,0}\definecolor{BLUE}{rgb}{0,0,1} %DIF PREAMBLE
\providecommand{\DIFadd}[1]{{\protect\color{blue}\uwave{#1}}} %DIF PREAMBLE
\providecommand{\DIFdel}[1]{{\protect\color{red}\sout{#1}}}                      %DIF PREAMBLE
%DIF SAFE PREAMBLE %DIF PREAMBLE
\providecommand{\DIFaddbegin}{} %DIF PREAMBLE
\providecommand{\DIFaddend}{} %DIF PREAMBLE
\providecommand{\DIFdelbegin}{} %DIF PREAMBLE
\providecommand{\DIFdelend}{} %DIF PREAMBLE
\providecommand{\DIFmodbegin}{} %DIF PREAMBLE
\providecommand{\DIFmodend}{} %DIF PREAMBLE
%DIF FLOATSAFE PREAMBLE %DIF PREAMBLE
\providecommand{\DIFaddFL}[1]{\DIFadd{#1}} %DIF PREAMBLE
\providecommand{\DIFdelFL}[1]{\DIFdel{#1}} %DIF PREAMBLE
\providecommand{\DIFaddbeginFL}{} %DIF PREAMBLE
\providecommand{\DIFaddendFL}{} %DIF PREAMBLE
\providecommand{\DIFdelbeginFL}{} %DIF PREAMBLE
\providecommand{\DIFdelendFL}{} %DIF PREAMBLE
\newcommand{\DIFscaledelfig}{0.5}
%DIF HIGHLIGHTGRAPHICS PREAMBLE %DIF PREAMBLE
\RequirePackage{settobox} %DIF PREAMBLE
\RequirePackage{letltxmacro} %DIF PREAMBLE
\newsavebox{\DIFdelgraphicsbox} %DIF PREAMBLE
\newlength{\DIFdelgraphicswidth} %DIF PREAMBLE
\newlength{\DIFdelgraphicsheight} %DIF PREAMBLE
% store original definition of \includegraphics %DIF PREAMBLE
\LetLtxMacro{\DIFOincludegraphics}{\includegraphics} %DIF PREAMBLE
\newcommand{\DIFaddincludegraphics}[2][]{{\color{blue}\fbox{\DIFOincludegraphics[#1]{#2}}}} %DIF PREAMBLE
\newcommand{\DIFdelincludegraphics}[2][]{% %DIF PREAMBLE
\sbox{\DIFdelgraphicsbox}{\DIFOincludegraphics[#1]{#2}}% %DIF PREAMBLE
\settoboxwidth{\DIFdelgraphicswidth}{\DIFdelgraphicsbox} %DIF PREAMBLE
\settoboxtotalheight{\DIFdelgraphicsheight}{\DIFdelgraphicsbox} %DIF PREAMBLE
\scalebox{\DIFscaledelfig}{% %DIF PREAMBLE
	\parbox[b]{\DIFdelgraphicswidth}{\usebox{\DIFdelgraphicsbox}\\[-\baselineskip] \rule{\DIFdelgraphicswidth}{0em}}\llap{\resizebox{\DIFdelgraphicswidth}{\DIFdelgraphicsheight}{% %DIF PREAMBLE
			\setlength{\unitlength}{\DIFdelgraphicswidth}% %DIF PREAMBLE
			\begin{picture}(1,1)% %DIF PREAMBLE
			\thicklines\linethickness{2pt} %DIF PREAMBLE
			{\color[rgb]{1,0,0}\put(0,0){\framebox(1,1){}}}% %DIF PREAMBLE
			{\color[rgb]{1,0,0}\put(0,0){\line( 1,1){1}}}% %DIF PREAMBLE
			{\color[rgb]{1,0,0}\put(0,1){\line(1,-1){1}}}% %DIF PREAMBLE
			\end{picture}% %DIF PREAMBLE
		}\hspace*{3pt}}} %DIF PREAMBLE
} %DIF PREAMBLE
\LetLtxMacro{\DIFOaddbegin}{\DIFaddbegin} %DIF PREAMBLE
\LetLtxMacro{\DIFOaddend}{\DIFaddend} %DIF PREAMBLE
\LetLtxMacro{\DIFOdelbegin}{\DIFdelbegin} %DIF PREAMBLE
\LetLtxMacro{\DIFOdelend}{\DIFdelend} %DIF PREAMBLE
\DeclareRobustCommand{\DIFaddbegin}{\DIFOaddbegin \let\includegraphics\DIFaddincludegraphics} %DIF PREAMBLE
\DeclareRobustCommand{\DIFaddend}{\DIFOaddend \let\includegraphics\DIFOincludegraphics} %DIF PREAMBLE
\DeclareRobustCommand{\DIFdelbegin}{\DIFOdelbegin \let\includegraphics\DIFdelincludegraphics} %DIF PREAMBLE
\DeclareRobustCommand{\DIFdelend}{\DIFOaddend \let\includegraphics\DIFOincludegraphics} %DIF PREAMBLE
\LetLtxMacro{\DIFOaddbeginFL}{\DIFaddbeginFL} %DIF PREAMBLE
\LetLtxMacro{\DIFOaddendFL}{\DIFaddendFL} %DIF PREAMBLE
\LetLtxMacro{\DIFOdelbeginFL}{\DIFdelbeginFL} %DIF PREAMBLE
\LetLtxMacro{\DIFOdelendFL}{\DIFdelendFL} %DIF PREAMBLE
\DeclareRobustCommand{\DIFaddbeginFL}{\DIFOaddbeginFL \let\includegraphics\DIFaddincludegraphics} %DIF PREAMBLE
\DeclareRobustCommand{\DIFaddendFL}{\DIFOaddendFL \let\includegraphics\DIFOincludegraphics} %DIF PREAMBLE
\DeclareRobustCommand{\DIFdelbeginFL}{\DIFOdelbeginFL \let\includegraphics\DIFdelincludegraphics} %DIF PREAMBLE
\DeclareRobustCommand{\DIFdelendFL}{\DIFOaddendFL \let\includegraphics\DIFOincludegraphics} %DIF PREAMBLE
%DIF LISTINGS PREAMBLE %DIF PREAMBLE
\RequirePackage{listings} %DIF PREAMBLE
\RequirePackage{color} %DIF PREAMBLE
\lstdefinelanguage{DIFcode}{ %DIF PREAMBLE
%DIF DIFCODE_UNDERLINE %DIF PREAMBLE
moredelim=[il][\color{red}\sout]{\%DIF\ <\ }, %DIF PREAMBLE
moredelim=[il][\color{blue}\uwave]{\%DIF\ >\ } %DIF PREAMBLE
} %DIF PREAMBLE
\lstdefinestyle{DIFverbatimstyle}{ %DIF PREAMBLE
language=DIFcode, %DIF PREAMBLE
basicstyle=\ttfamily, %DIF PREAMBLE
columns=fullflexible, %DIF PREAMBLE
keepspaces=true %DIF PREAMBLE
} %DIF PREAMBLE
\lstnewenvironment{DIFverbatim}{\lstset{style=DIFverbatimstyle}}{} %DIF PREAMBLE
\lstnewenvironment{DIFverbatim*}{\lstset{style=DIFverbatimstyle,showspaces=true}}{} %DIF PREAMBLE
%DIF END PREAMBLE EXTENSION ADDED BY LATEXDIFF
%%

%Title, date an author of the document
%\title{Bibliography management: BibTeX}
%\author{Overleaf}

%Begining of the document
\begin{document}

%%% AAB - Mudei o titulo
\title{Fast Noise Reduction in InSAR Data with Empirical Statistical Models\\
	Revision R1}


\author{Anderson A.\ De Borba,
	Joselito E.\ De Araújo, and
	Alejandro C. Frery}


\maketitle


\section{Review Report}
\begin{tcolorbox}[colback=red!5!white,colframe=red!75!black,title=Review Report]
This manuscript addresses an interesting topic in InSAR data processing, but Major revisions
are required to improve clarity, completeness, and experimental rigor before publication.

1. The descriptions of models in abstract and highlights may be confusing. For example, “two
empirical models”, “a model based on the characteristics of sensing”, “restricted version of the
physical model”, “two physical models” and “two empirical circular models”.

2. Please provide the full names of the filters before using them, “LInSARREF”, “TcNfilter” and
“TcCfilter”.

3. Please provide more specific details about the physical models and empirical models in the
introduction.

4. Please specify the contributions of this paper in the introduction.

5. Between lines 273–276: Are the phase limits identical for all pixels? Since coherence is also
related to scattering type, the authors should provide comparisons for different land cover types.

6. Figure 9: Please supplement the figure with an analysis across different land cover types.
\end{tcolorbox}

%%%%%%%%%%%%%%%%%%%%%%%%%%%%%%%%%%%%%%%%%%%%%%%%%%%%%%%%%%%%%%%%%%%%%%%%%%%%%%%%%%
%%%%%%%%%%%%%%%%%%%%%%%%%%% Comment #Letter %%%%%%%%%%%%%%%%%%%%%%%%%%%%%%%%%%%%%%%%%%%
%%%%%%%%%%%%%%%%%%%%%%%%%%%%%%%%%%%%%%%%%%%%%%%%%%%%%%%%%%%%%%%%%%%%%%%%%%%%%%%%%%


\section{Answers to Reviewer. \#2}

Thank you very much for handling this manuscript.

We have prepared a revised version taking into account all the comments and suggestions made.

In fact, we found the review well-informed and constructive. We would like to thank the reviewer.

This response letter addresses all the comments in red, followed by
our reactions, and, whenever necessary, the changes made.

The authors
%%%%%%%%%%%%%%%%%%%%%%%%%%%%%%%%%%%%%%%%%%%%%%%%%%%%%%%%%%%%%%%%%%%%%%%%%%%%%%%%%%
%%%%%%%%%%%%%%%%%%%%%%%%%%% Comment #1 %%%%%%%%%%%%%%%%%%%%%%%%%%%%%%%%%%%%%%%%%%%
%%%%%%%%%%%%%%%%%%%%%%%%%%%%%%%%%%%%%%%%%%%%%%%%%%%%%%%%%%%%%%%%%%%%%%%%%%%%%%%%%%
\vskip3em\begin{tcolorbox}[colback=red!5!white,colframe=red!75!black,title=Comment 1]
% AAB inserido
1. The descriptions of models in abstract and highlights may be confusing. For example, “two
empirical models”, “a model based on the characteristics of sensing”, “restricted version of the
physical model”, “two physical models” and “two empirical circular models”.
\end{tcolorbox}

\begin{tcolorbox}[colback=white,colframe=black,title=Changes 1]
Thank you very much for your contribution. The authors fully agree
with the suggestion, and we have made the following modification in the Abstract and Highlights:

- Abstract:

\DIFdelbegin \DIFdel{However, the noise inherent in interferometric images
reduces the accuracy and reliability of these products. The development of new methods
and techniques for filtering InSAR images has attracted the attention of several researchers
seeking to reduce noise while preserving image details, resolution, and structures. This
article analyzes the impact of using two empirical models for noise reduction, as well
as a model based on the characteristics of sensing with coherent illumination.
}\DIFdelend \DIFaddbegin \DIFadd{However, the inherent noise in interferometric images significantly compromises the accuracy and reliability of the resulting data products. Consequently, the development of advanced InSAR filtering techniques has become a focal point for researchers aiming to suppress noise while simultaneously preserving fine details, spatial resolution, and structural integrity. This article evaluates the performance of three distinct filtering approaches: a model based on the physical characteristics of sensing with coherent illumination, known as the Gierull model, and two empirical models, the truncated wrapped normal (TcN) and truncated wrapped Cauchy (TcC) distributions.
}\DIFaddend %
\DIFdelbegin \DIFdel{These
empirical models drastically reduce the processing time without compromising the results.
Additionally, we quantify the effect of using a restricted version of the physical model. To
}\DIFdelend \DIFaddbegin \DIFadd{These models drastically reduce the processing time without compromising the results. To
}\DIFaddend confirm this affirmation, we use four types of test images, including a simulated dataset for
goal validation of the proposed methods. We use La Cumbre volcano, Los Alamos, and
Robledo volcano InSAR images to evaluate performance on real InSAR images.
%
The study found a bottleneck in the computational cost of numerical integration\DIFdelbegin \DIFdel{ that}\DIFdelend \DIFaddbegin \DIFadd{, which}\DIFaddend{} the empirical
models impose on the filtering, a topic that requires further investigation.
%
In this article\DIFdelbegin \DIFdel{,
the authors propose}\DIFdelend \DIFaddbegin \DIFadd{, the authors propose}\DIFaddend{} three new techniques that demonstrate great computational efficiency
compared to the cutting-edge methods for filtering InSAR images with equivalent accuracy.

- Highlights:

\begin{itemize}[label=$\bullet$]
\item We quantified the \DIFdelbegin \DIFdel{difference between two physical }\DIFdelend \DIFaddbegin \DIFadd{differences between the Lee and Gierull }\DIFaddend models (a complete model and a \DIFdelbegin \DIFdel{mode with }\DIFdelend \DIFaddbegin \DIFadd{model with a }\DIFaddend restricted parameter space) for InSAR data.
 \DIFaddbegin 

\DIFaddend \item We proposed two \DIFdelbegin \DIFdel{empirical circular models }\DIFdelend \DIFaddbegin \DIFadd{new filters }\DIFaddend for noise reduction \DIFdelbegin \DIFdel{.
    }\DIFdelend \DIFaddbegin \DIFadd{based on the truncated wrapped normal (TcN) and truncated wrapped Cauchy (TcC) models.
}

\DIFaddend \item We show that these empirical models produce \DIFaddbegin \DIFadd{results as good }\DIFaddend as \DIFdelbegin \DIFdel{good results as }\DIFdelend the physical one\DIFaddbegin \DIFadd{, }\DIFaddend with a significant reduction in computational \DIFdelbegin \DIFdel{processing }\DIFdelend time.
\end{itemize}
\end{tcolorbox}


%%%%%%%%%%%%%%%%%%%%%%%%%%%%%%%%%%%%%%%%%%%%%%%%%%%%%%%%%%%%%%%%%%%%%%%%%%%%%%%%%%
%%%%%%%%%%%%%%%%%%%%%%%%%%% Comment #2 %%%%%%%%%%%%%%%%%%%%%%%%%%%%%%%%%%%%%%%%%%%
%%%%%%%%%%%%%%%%%%%%%%%%%%%%%%%%%%%%%%%%%%%%%%%%%%%%%%%%%%%%%%%%%%%%%%%%%%%%%%%%%%
\vskip3em\begin{tcolorbox}[colback=red!5!white,colframe=red!75!black,title=Comment 2]
% AAB inserido
2. Please provide the full names of the filters before using them, “LInSARREF”, “TcNfilter” and
“TcCfilter”.
\end{tcolorbox}

\begin{tcolorbox}[colback=white,colframe=black,title=Changes 1]
Thank you very much for your contribution. The authors fully agree
with the suggestion, and we have made the following modification in the Introduction:

To address the \DIFdelbegin \DIFdel{need }\DIFdelend \DIFaddbegin \DIFadd{demand }\DIFaddend for fast and effective interferometric \DIFdelbegin \DIFdel{phase‑filtering }\DIFdelend \DIFaddbegin \DIFadd{phase-filtering }\DIFaddend techniques, we propose three \DIFdelbegin \DIFdel{new filters: LInSARRFE, TcNfilter, and TcCfilter.
LInSARRFE }\DIFdelend \DIFaddbegin \DIFadd{novel filters: a filter based on the Gierull model (LInSARRFE), a filter based on the truncated wrapped normal model (TcNfilter), and a filter based on the truncated wrapped Cauchy model (TcCfilter). The LInSARRFE filter }\DIFaddend is derived from a physically based probability distribution that models SAR signal formation and its representation as a phase image\DIFdelbegin \DIFdel{.
}\DIFdelend \DIFaddbegin \DIFadd{, specifically following Gierull Model. In contrast, the }\DIFaddend TcNfilter and TcCfilter \DIFdelbegin \DIFdel{, in contrast, }\DIFdelend rely on empirical probability distributions designed to \DIFdelbegin \DIFdel{reduce computational complexity.
}\DIFdelend \DIFaddbegin \DIFadd{minimize computational complexity; these are based on the truncated wrapped normal (TcN) and truncated wrapped Cauchy distributions (TcC), respectively. 
}

\end{tcolorbox}


%%%%%%%%%%%%%%%%%%%%%%%%%%%%%%%%%%%%%%%%%%%%%%%%%%%%%%%%%%%%%%%%%%%%%%%%%%%%%%%%%%
%%%%%%%%%%%%%%%%%%%%%%%%%%% Comment #3 %%%%%%%%%%%%%%%%%%%%%%%%%%%%%%%%%%%%%%%%%%%
%%%%%%%%%%%%%%%%%%%%%%%%%%%%%%%%%%%%%%%%%%%%%%%%%%%%%%%%%%%%%%%%%%%%%%%%%%%%%%%%%%
\vskip3em\begin{tcolorbox}[colback=red!5!white,colframe=red!75!black,title=Comment 3]
% AAB inserido
3. Please provide more specific details about the physical models and empirical models in the
introduction.
\end{tcolorbox}

\begin{tcolorbox}[colback=white,colframe=black,title=Changes 3]
Thank you very much for your contribution. The authors fully agree
with the suggestion, and we have made the following modification in the Introduction:

\DIFadd{One important aspect of filtering methods is the distribution model used for the probability density function (PDF). In a series of articles by Lee~\mbox{%DIFAUXCMD
\cite{ImprovedSigmaFilterforSpeckleFilteringofSARImagery, AnewtechniquefornoisefilteringofSARinterferometricphasimages, RefinedFilteringofInterferometricPhaseFromInSARData}}\hskip0pt%DIFAUXCMD
, the Lee Filter and its variants were proposed, based on a model grounded in the physical properties of SAR image theory. The model was proposed in the article by~\mbox{%DIFAUXCMD
\cite{IntensityAndPhaseStatisticsOfMultilookPolarimetricAndInterferometricSARimagery}}\hskip0pt%DIFAUXCMD
. Gierrull, in the article, proposed a similar model that we used to construct the filter based on the problem's physics, as introduced in this article. On the other hand, given the periodic nature of the phase data in InSAR images, we proposed using empirical circular models to replace the Lee model in the filtering methods described. The circular models are described in~\mbox{%DIFAUXCMD
\cite{CircularStatisticsinR}}\hskip0pt%DIFAUXCMD
.
}


\end{tcolorbox}

%%%%%%%%%%%%%%%%%%%%%%%%%%%%%%%%%%%%%%%%%%%%%%%%%%%%%%%%%%%%%%%%%%%%%%%%%%%%%%%%%%
%%%%%%%%%%%%%%%%%%%%%%%%%%% Comment #4 %%%%%%%%%%%%%%%%%%%%%%%%%%%%%%%%%%%%%%%%%%%
%%%%%%%%%%%%%%%%%%%%%%%%%%%%%%%%%%%%%%%%%%%%%%%%%%%%%%%%%%%%%%%%%%%%%%%%%%%%%%%%%%
\vskip3em\begin{tcolorbox}[colback=red!5!white,colframe=red!75!black,title=Comment 4]
% AAB inserido
4. Please specify the contributions of this paper in the introduction.
\end{tcolorbox}

\begin{tcolorbox}[colback=white,colframe=black,title=Changes 4]
Thank you very much for your contribution. The authors fully agree
with the suggestion, and we have made the following modification in the Introduction:


In this work, we \DIFdelbegin \DIFdel{quantify the differences between Lee et al's and }\DIFdelend \DIFaddbegin \DIFadd{propose a novel methodology to reduce computational cost and improve the numerical stability of filters applied to InSAR images. Three models are introduced to modify conventional filtering approaches: a physics-based formulation, }\DIFaddend Gierull's \DIFdelbegin \DIFdel{versions of the same physical }\DIFdelend \DIFaddbegin \DIFadd{model, derived from the Lee }\DIFaddend model, and \DIFdelbegin \DIFdel{introduce two empirical distributions designed to further reduce computational cost and }\DIFdelend \DIFaddbegin \DIFadd{two additional empirical models. All three models are designed to reduce computational demands further while enhancing }\DIFaddend numerical stability. \DIFaddbegin \DIFadd{The core contribution of the proposed methodology lies in replacing the standard Lee model with either Gierull's model or one of the two empirical alternatives. }\DIFaddend These models are incorporated into an adaptive filtering framework that adjusts to local phase variability. \DIFdelbegin \DIFdel{We used the }\DIFdelend \DIFaddbegin \DIFadd{The }\DIFaddend refined Lee filter \DIFaddbegin \DIFadd{is adopted }\DIFaddend as a reference \DIFdelbegin \DIFdel{, as it represents the state of the art }\DIFdelend \DIFaddbegin \DIFadd{baseline, given its established role as a benchmark }\DIFaddend in InSAR data filtering.



\end{tcolorbox}

%%%%%%%%%%%%%%%%%%%%%%%%%%%%%%%%%%%%%%%%%%%%%%%%%%%%%%%%%%%%%%%%%%%%%%%%%%%%%%%%%%
%%%%%%%%%%%%%%%%%%%%%%%%%%% Comment #5 %%%%%%%%%%%%%%%%%%%%%%%%%%%%%%%%%%%%%%%%%%%
%%%%%%%%%%%%%%%%%%%%%%%%%%%%%%%%%%%%%%%%%%%%%%%%%%%%%%%%%%%%%%%%%%%%%%%%%%%%%%%%%%
\vskip3em\begin{tcolorbox}[colback=red!5!white,colframe=red!75!black,title=Comment 5]
% AAB inserido
Between lines 273–276: Are the phase limits identical for all pixels? Since coherence is also
related to scattering type, the authors should provide comparisons for different land cover types.
\end{tcolorbox}

\begin{tcolorbox}[colback=white,colframe=black,title=Changes 5]

\end{tcolorbox}

%%%%%%%%%%%%%%%%%%%%%%%%%%%%%%%%%%%%%%%%%%%%%%%%%%%%%%%%%%%%%%%%%%%%%%%%%%%%%%%%%%
%%%%%%%%%%%%%%%%%%%%%%%%%%% Comment #6 %%%%%%%%%%%%%%%%%%%%%%%%%%%%%%%%%%%%%%%%%%%
%%%%%%%%%%%%%%%%%%%%%%%%%%%%%%%%%%%%%%%%%%%%%%%%%%%%%%%%%%%%%%%%%%%%%%%%%%%%%%%%%%
\vskip3em\begin{tcolorbox}[colback=red!5!white,colframe=red!75!black,title=Comment 6]
% AAB inserido
6. Figure 9: Please supplement the figure with an analysis across different land cover types.
\end{tcolorbox}

\begin{tcolorbox}[colback=white,colframe=black,title=Changes 6]

\end{tcolorbox}

\bibliographystyle{abbrvnat}

%Sets the bibliography style to UNSRT and imports the 
%bibliography file "sample.bib".
\bibliography{references}
\end{document}